\chapter{Học tăng cường theo Naeem và cộng sự (2020)}\label{chap:2}

Bài tổng quan \emph{A Gentle Introduction to Reinforcement Learning and its Application in Different Fields} của
Naeem, Rizvi và Coronato (IEEE Access, tập 8, tr.~209320--209344)~\cite{ch2:naeem} hệ thống hóa học tăng cường
(reinforcement learning, RL) theo ba lớp: khung toán học MDP (phần II--III), các họ thuật toán (phần IV) và ứng
dụng cùng thách thức (phần V). Chương này đi theo đúng trình tự đó, đối chiếu với giáo trình của Sutton và
Barto~\cite{ch2:sutton}, và chỉ ra phần nào được dùng trực tiếp trong thí nghiệm ở Chương~\ref{chap:4}--\ref{chap:7}.

\section{Quá trình quyết định Markov}

\begin{dinhnghia}{Định nghĩa 2.1 (MDP)}
Một quá trình quyết định Markov (Markov decision process) là bộ $\mathcal M=(\mathcal S,\mathcal A,P,R,\gamma)$
gồm tập trạng thái $\mathcal S$, tập hành động $\mathcal A$, xác suất chuyển $P(s'\mid s,a)$, hàm phần thưởng
$R(s,a)$ và hệ số chiết khấu $\gamma\in[0,1)$. Tính Markov yêu cầu
$\Pr(s_{t+1},r_t\mid s_t,a_t,s_{t-1},\dots)=\Pr(s_{t+1},r_t\mid s_t,a_t)$.
\end{dinhnghia}

Tác tử và môi trường tương tác theo các bước rời rạc (Hình~\ref{fig:ch2-loop}): ở bước $t$ tác tử thấy $s_t$, chọn
$a_t\sim\pi(\cdot\mid s_t)$, nhận phần thưởng $r_t$ và trạng thái mới $s_{t+1}$. Mục tiêu không phải phần thưởng
tức thời mà là \textbf{lợi tức} (return) chiết khấu, và chất lượng của chính sách được đo bằng hàm giá trị:
\begin{gather}
G_t=\sum_{k=0}^{T-t-1}\gamma^{k}r_{t+k},\label{eq:ch2-return}\\
V^{\pi}(s)=\mathbb E_{\pi}\!\left[G_t\mid s_t=s\right],\qquad
Q^{\pi}(s,a)=\mathbb E_{\pi}\!\left[G_t\mid s_t=s,a_t=a\right].\notag
\end{gather}

\hinhtikz{ch2_loop}{Vòng tương tác tác tử -- môi trường của MDP và quỹ đạo $(s_0,a_0,r_0,s_1,\dots)$ sinh ra theo thời gian.}{ch2-loop}

\nhanxet{Phần thưởng $r_t$ là phản hồi một bước, còn $G_t$ cộng dồn có chiết khấu. Trong bài thí nghiệm, báo cáo
lợi nhuận \emph{trung bình không chiết khấu} mỗi bước, trong khi bảng Q ước lượng lợi tức với $\gamma=0{,}9$; hai đại
lượng vì vậy khác thang đo dù cùng đơn vị tiền.}

\section{Phương trình Bellman}

Tách số hạng đầu khỏi tổng~\eqref{eq:ch2-return} cho quan hệ đệ quy giữa giá trị hiện tại và giá trị ở bước sau.
Với một chính sách cố định $\pi$ ta có \textbf{phương trình Bellman kỳ vọng}, còn với chính sách tối ưu ta có
\textbf{phương trình Bellman tối ưu}~\cite{ch2:naeem,ch2:sutton}:
\begin{align}
V^{\pi}(s)&=\sum_{a}\pi(a\mid s)\sum_{s'}P(s'\mid s,a)\big[R(s,a)+\gamma V^{\pi}(s')\big],\label{eq:ch2-bell-v}\\
Q^{*}(s,a)&=\sum_{s'}P(s'\mid s,a)\Big[R(s,a)+\gamma\max_{a'}Q^{*}(s',a')\Big],\label{eq:ch2-bell-q}\\
\pi^{*}(s)&=\arg\max_{a}Q^{*}(s,a).\notag
\end{align}
Khi biết $P$ và $R$, \eqref{eq:ch2-bell-q} giải được bằng lặp giá trị; khi không biết, các thuật toán RL thay kỳ vọng
bằng mẫu trải nghiệm. Bài thí nghiệm nhỏ đến mức mô hình $P,R$ tính được chính xác, nên Chương~\ref{chap:4} dùng
\eqref{eq:ch2-bell-q} làm giá trị đối chiếu cho những gì Q-learning học được.

\section{Các họ thuật toán}

Naeem và cộng sự chia thuật toán theo hai trục: có dùng mô hình môi trường hay không (model-based /
model-free), và học giá trị hay học trực tiếp chính sách (Hình~\ref{fig:ch2-tax}).

\hinhtikz{ch2_tax}{Phân loại các thuật toán RL theo bài tổng quan; ô viền đỏ là thuật toán được dùng trong thí nghiệm.}{ch2-tax}

\begin{itemize}
\item \textbf{Quy hoạch động} (dynamic programming, DP): lặp chính sách / lặp giá trị trên mô hình đã biết; chính xác
nhưng cần $P$ và tốn kém khi $|\mathcal S|$ lớn.
\item \textbf{Monte Carlo} (MC): ước lượng $V^\pi$ bằng trung bình lợi tức của các episode hoàn chỉnh; không
\emph{bootstrap}, không thiên lệch nhưng phương sai lớn và phải chờ hết episode.
\item \textbf{Sai phân thời gian} (temporal difference, TD): cập nhật sau từng bước bằng ước lượng của chính nó,
$V(s_t)\leftarrow V(s_t)+\alpha\,[r_t+\gamma V(s_{t+1})-V(s_t)]$; kết hợp lấy mẫu của MC với bootstrap của DP.
\item \textbf{Actor--Critic}, phương pháp Bayes và \textbf{DQN}: tách chính sách (actor) khỏi giá trị (critic), mô hình
hóa bất định, hoặc xấp xỉ $Q$ bằng mạng nơ-ron cho không gian lớn.
\end{itemize}

\subsection{Q-learning và SARSA}

Hai đại diện quan trọng nhất của TD điều khiển chỉ khác nhau ở \emph{mục tiêu} (target) của phép cập nhật
(Hình~\ref{fig:ch2-backup}). Q-learning~\cite{ch2:watkins} là \textbf{off-policy}: mục tiêu lấy hành động tốt nhất
ở trạng thái sau, bất kể hành động thực sự được chọn; SARSA là \textbf{on-policy}: mục tiêu dùng chính hành động
$a_{t+1}$ mà chính sách đang khám phá chọn:
\begin{align}
\text{Q-learning:}&\quad Q(s_t,a_t)\leftarrow Q(s_t,a_t)+\alpha\Big[\underbrace{r_t+\gamma\max_{a'}Q(s_{t+1},a')}_{y_t}-Q(s_t,a_t)\Big],\label{eq:ch2-q}\\
\text{SARSA:}&\quad Q(s_t,a_t)\leftarrow Q(s_t,a_t)+\alpha\big[r_t+\gamma\,Q(s_{t+1},a_{t+1})-Q(s_t,a_t)\big].\label{eq:ch2-sarsa}
\end{align}
Ở bước cuối episode, mục tiêu chỉ còn $y_t=r_t$. Đại lượng $\delta_t=y_t-Q(s_t,a_t)$ là \textbf{sai số TD}.

\hinhtikz{ch2_backup}{Sơ đồ sao lưu (backup diagram) của Q-learning (lấy max) và SARSA (lấy hành động thực).}{ch2-backup}

\nhanxet{Vì Q-learning học giá trị của chính sách tham lam trong khi vẫn khám phá, nó cho phép tách gọn pha huấn
luyện ($\varepsilon>0$) khỏi pha đánh giá ($\varepsilon=0$) --- đúng cách đánh giá được dùng ở Chương~\ref{chap:5}.}

\subsection{Khám phá, tốc độ học và điều kiện hội tụ}

Chiến lược \textbf{$\varepsilon$-greedy} chọn ngẫu nhiên với xác suất $\varepsilon$, ngược lại chọn
$\arg\max_a Q(s,a)$; với hai hành động, hành động ``không tốt nhất'' được chọn với xác suất $\varepsilon/2$.
Thường $\varepsilon$ giảm dần, ví dụ $\varepsilon_e=\max(\varepsilon_{\min},\varepsilon_0\rho^{e})$ theo episode $e$.
Tốc độ học $\alpha$ quyết định mỗi mẫu thay đổi $Q$ bao nhiêu: $\alpha$ nhỏ cho đường học mượt nhưng chậm, $\alpha$
lớn phản ứng nhanh nhưng khuếch đại nhiễu. Q-learning hội tụ về $Q^{*}$ với xác suất 1 nếu MDP dừng, mọi cặp
$(s,a)$ được thăm vô hạn lần và $\sum_t\alpha_t=\infty$, $\sum_t\alpha_t^2<\infty$~\cite{ch2:watkins}.

\begin{luuy}{Khoảng cách giữa lý thuyết và thí nghiệm}
Thí nghiệm dùng $\alpha$ hằng, số bước hữu hạn và một trạng thái rút gọn không chứa hành động trước của chính tác tử;
đối thủ cũng có thể học. Cả ba điều kiện trên vì vậy không được bảo đảm hoàn toàn, nên độ ổn định phải được đo bằng
nhiều seed thay vì suy ra từ định lý.
\end{luuy}

\subsection{Từ bảng Q tới DQN}

Khi $|\mathcal S|$ lớn, bảng Q được thay bằng mạng $Q(s,a;\theta)$. DQN~\cite{ch2:mnih} ổn định quá trình học bằng
hai cơ chế: \emph{bộ nhớ phát lại} (experience replay) phá tương quan giữa các mẫu liên tiếp, và \emph{mạng mục tiêu}
$\theta^{-}$ cập nhật chậm; hàm mất mát là
$\mathcal L(\theta)=\mathbb E\big[(r+\gamma\max_{a'}Q(s',a';\theta^{-})-Q(s,a;\theta))^2\big]$. Bài thí nghiệm chỉ có
$4\times2=8$ giá trị Q nên bảng là lựa chọn đúng; DQN được nêu để chỉ hướng mở rộng khi trạng thái giàu hơn.

\section{Ứng dụng và thách thức}\label{sec:ch2-challenge}

Phần V của bài tổng quan khảo sát RL trong nhiều lĩnh vực; điểm chung là chuỗi quyết định nối tiếp, trong đó
hành động hôm nay thay đổi tình huống ngày mai (Bảng~\ref{tab:ch2-app}).

\begin{table}[!htbp]\centering\caption{Một số lĩnh vực ứng dụng RL được Naeem và cộng sự khảo sát.}\label{tab:ch2-app}
{\small\begin{tabular}{@{}L{3.6cm}L{5.4cm}L{5.5cm}@{}}\toprule
Lĩnh vực & Trạng thái / hành động & Phần thưởng điển hình \\\midrule
Robot, điều khiển & Cảm biến khớp / mô-men xoắn & Bám quỹ đạo, tiết kiệm năng lượng \\
Mạng, viễn thông & Tải kênh / cấp phát tài nguyên & Thông lượng, độ trễ \\
IoT tự chủ & Trạng thái thiết bị / lịch hoạt động & Tuổi thọ pin, chất lượng dịch vụ \\
Ngôn ngữ, thị giác & Ngữ cảnh / từ hoặc vùng chọn & Điểm chất lượng đầu ra \\
Trò chơi & Bàn cờ, khung hình / nước đi & Thắng / thua, điểm \\
Lập lịch, tài nguyên & Hàng đợi / phân công & Thời gian hoàn thành \\
Y tế & Tình trạng bệnh nhân / phác đồ & Kết cục điều trị \\\bottomrule
\end{tabular}}
\end{table}

Các thách thức được nêu gồm: \textbf{hiệu quả mẫu} (cần rất nhiều tương tác), \textbf{không gian lớn},
\textbf{an toàn} khi khám phá trên hệ thật, \textbf{quan sát một phần}, \textbf{môi trường phi dừng} và đánh đổi
\textbf{thiên lệch--phương sai}. Cả sáu đều xuất hiện ở quy mô nhỏ trong bài thí nghiệm: doanh số nhị thức tạo phương
sai; bootstrap tạo thiên lệch; trạng thái thiếu hành động trước của chính mình làm quan sát không đầy đủ; nếu đối thủ
cũng học, phân phối phản hồi thay đổi theo thời gian; và mỗi lần ``khám phá'' bằng giảm giá là một khoản lợi nhuận
thật bị mất. Vì vậy thí nghiệm lấy từ bài tổng quan bốn thành phần: MDP (mô hình hóa), Q-learning off-policy với
$\varepsilon$-greedy giảm dần (tác tử A), Bellman tối ưu trên mô hình đã biết (giá trị đối chiếu), và thiết kế
đánh giá nhiều seed cùng thí nghiệm hai tác tử cùng học (đối phó phương sai và tính phi dừng).
